\documentclass[11pt,a4paper]{article}
\usepackage[T1]{fontenc}
\usepackage[utf8]{inputenc}
\usepackage{lmodern}
\usepackage[provide=*,polish]{babel}
\usepackage[margin=2.15cm]{geometry}
\usepackage[hidelinks]{hyperref}
\pagestyle{empty}
\setlength{\parindent}{0pt}
\setlength{\parskip}{4pt}
\newcommand{\dzial}[1]{\par\vspace{8pt}{\large\bfseries #1}\par\vspace{1pt}\hrule\vspace{4pt}}

\begin{document}
{\LARGE\bfseries Mateusz Nawrot}\par
Programowanie w C++\par
\href{https://github.com/nawrotmateusz56-sudo}{github.com/nawrotmateusz56-sudo}

\dzial{Profil}
W C++ zajmuję się algorytmami i programowaniem obiektowym. Lubię zadania,
w których wybór struktury danych wpływa na prostotę rozwiązania.
Piszę też w Pythonie, C i C\#.

\dzial{Wykształcenie}
\textbf{Technikum informatyczne}

\dzial{Programowanie w C++}
\textbf{Organizacja kodu.}
Dzielę rozwiązanie na funkcje o określonym zadaniu. Gdy dane i operacje
na nich tworzą jedną całość, porządkuję je w klasie. Dzięki takiemu podziałowi
mogę sprawdzić pojedynczą funkcję niezależnie od reszty programu.\par
\textbf{Algorytmy i struktury danych.}
Dobieram sposób przechowywania danych do potrzeb programu. Przy pracy
z kolekcją biorę pod uwagę, czy częściej będę przechodzić po jej elementach,
wyszukiwać wartości, czy wstawiać nowe dane.\par
\textbf{Testowanie i debugowanie.}
Sprawdzam wynik dla typowych danych i przypadków brzegowych, takich jak
pusty zbiór, jeden element i powtarzające się wartości. Przy błędzie
odtwarzam go na małym przykładzie i śledzę wartości zmiennych w debuggerze.

\dzial{Narzędzia i pozostałe umiejętności}
\textbf{Środowisko pracy:} Visual Studio, Visual Studio Code i terminal.\par
\textbf{Kontrola wersji:} Git i GitHub.\par
\textbf{Bazy danych:} SQL.\par
\textbf{Testy:} pisanie testów jednostkowych.

\dzial{Języki obce}
Angielski — C2.
\end{document}
